\documentclass[UTF8,fontset=fandol,zihao=5]{ctexart}

\usepackage[a4paper,margin=2cm]{geometry}
\usepackage{amsmath,amssymb,mathtools}
\usepackage{upgreek}
\usepackage[nolimits]{cmupint}
\usepackage{booktabs,tabularx,array}
\usepackage{xcolor}
\usepackage[most]{tcolorbox}
\usepackage{enumitem}
\usepackage{indentfirst}
\usepackage{hyperref}
\usepackage{tikz}


\usetikzlibrary{
  positioning,
  calc,
  decorations.pathreplacing,
  arrows.meta
}

% relationbox 中知识关系图的统一样式
\tikzset{
  relnode/.style={
    font=\small,
    align=center
  },
  relline/.style={
    thick
  },
  relarrow/.style={
    -{Stealth},
    thick
  }
}

\DeclareTextFontCommand{\textbf}{\heiti\color{main}}

\definecolor{main}{HTML}{B33F46}
\definecolor{mainlight}{HTML}{FFF0F0}
\definecolor{method}{HTML}{B3797C}
\definecolor{methodlight}{HTML}{FFF0F0}

\hypersetup{colorlinks=true,linkcolor=main,urlcolor=main}
\setlength{\parindent}{2em}
\linespread{1.5}
\setlength{\abovedisplayskip}{5pt}
\setlength{\belowdisplayskip}{5pt}
\setlength{\abovedisplayshortskip}{3pt}
\setlength{\belowdisplayshortskip}{3pt}
\setlist[itemize]{leftmargin=2em,itemsep=0.1em,topsep=0.2em}
\setlist[enumerate]{leftmargin=2.3em,itemsep=0.1em,topsep=0.2em}

\ctexset{
  section={format=\centering\heiti\zihao{3}\color{main},
            name={第,章},number=\chinese{section},aftername=\quad,
            beforeskip=0.9em,afterskip=0.5em},
  subsection={format=\heiti\zihao{5}\color{main},
              name={,、},number=\chinese{subsection},aftername=,
              beforeskip=0.6em,afterskip=0.22em},
  subsubsection={format=\heiti\zihao{5}\color{main},
                  number={（\chinese{subsubsection}）},aftername=,
                  beforeskip=0.55em,afterskip=0.15em}
}

\newtcolorbox{relationbox}{
  enhanced,breakable,colback=mainlight,colframe=main,
  boxrule=0.65pt,arc=1.4mm,left=2mm,right=2mm,top=1.2mm,bottom=1.2mm
}
\newtcolorbox{methodbox}[1][]{
  enhanced,breakable,colback=methodlight,colframe=method,
  boxrule=0.65pt,arc=1.4mm,left=2mm,right=2mm,top=1.2mm,bottom=1.2mm,
  title=#1,fonttitle=\kaishu,fontupper=\kaishu,
  before upper={\setlength{\parindent}{2em}}
}

\newcommand{\ii}{\mathrm{i}}
\newcommand{\ee}{\mathrm{e}}
\newcommand{\CC}{\mathbb{C}}
\newcommand{\RR}{\mathbb{R}}
\newcommand{\dd}{\mathrm{d}}

\begin{document}

\begin{center}
  {\heiti\zihao{1}\color{main} 复分析笔记}
\end{center}

\begin{relationbox}
\[
  \boxed{\text{复数域}}
  \longrightarrow
  \boxed{\text{复函数}}
  \longrightarrow
  \boxed{\text{极限与连续}}
  \longrightarrow
  \boxed{\text{微分学}}
\]
\end{relationbox}

\section{复数域}
正如实分析的基础是Euclid空间$\RR^n$，复分析的基础是复数域$\CC$。复数域可视为实
二维向量空间 \(\operatorname{span}_{\mathbb R}\{1,\ii\}\) 上进一步赋予乘法结构
所得。规定基元素的乘法 \(1 \cdot 1=1,\ 1 \cdot \ii=\ii \cdot 1=\ii,\ \ii^2=-1\)，并将其按
双线性延拓至整个空间，即得到复数域。下面我们将详细描述如何从二维Euclid空
间$\RR^2$构造出复数域$\CC$。


\subsection{复数集}
先不加定义直接给出虚数单位$\ii$。由此，我们可以给出一组标准正交基$\left \{1, \ii\right \}$，
并张成二维Euclid空间$\RR^2$。故所有复数都可以对应到$\RR^2$上的点，并记为
$z=x \cdot 1 + y \cdot \ii$，与点$(x,y)$一一对应起来。由此，我们也不难理解复数
代数表示$z=x+\ii y$中的加号$+$是什么意思了，其表示的正是二维Euclid空间$\RR^2$
中的向量加法。

不难注意到，在加法、数乘、长度、角度这些层面上，复数域$\CC$与二维Euclid空间
$\RR^2$是完全一致的。也就是说作为实Euclid空间，复数域与二维Euclid空间同构，
即$\CC \cong \RR^2$。

在复数域$\CC$与二维Euclid空间$\RR^2$中，因为同构而定义相似的概念有很多，
现列举一些：
\begin{itemize}
  \item 向量/复数运算：加法、数乘、求长度/求模、求角度等；
  \item 集合论概念：开集、闭集等；
  \item 拓扑概念：邻域、边界、内点、聚点等；
  \item 几何概念：如复平面、Jordan曲线等；
  \item 分析概念：如极限、连续等。
  \item 复数间一般不能比较大小。
\end{itemize}

\subsection{复数域}

在复数集基础上进一步赋予乘法结构，即规定$1\cdot1=1$、$1\cdot \ii=\ii\cdot1=\ii$、
$\ii^2=-1$，
我们便得到了复数域$\CC$。这便是复数域$\CC$之区别于二维Euclid空间$\RR^2$的原因。
复数域$\CC$规定了复数间的乘法（注意不是数乘或叉乘），这使得复数域相对于二维
Euclid空间$\RR^2$有了更丰富的结构。

因乘法结构而区别于二维Euclid空间$\RR^2$的复数概念有：
\begin{itemize}
  \item 复数运算：乘除法、指对幂运算等；
  \item 分析概念：如导数、可微等。
\end{itemize}

除此以外，复数的旋转性质也是从此而来。注意到两向量基正交且长度为1，
故由$1 \cdot \ii = \ii$，$\ii \cdot \ii = -1$，可知道对任意向量乘以$\ii$，
两向量基均逆时针旋转了$\frac{\pi}{2}$，因此其结果是将该向量逆时针旋转了$\frac{\pi}{2}$。
这便是复数的旋转性质。

接下来给出复数的基本概念。

\begin{enumerate}
  \item 形式
  \begin{center}
    \begin{tabularx}{0.92\textwidth}{>{\bfseries}p{2.5cm}X X}
      \toprule
      形式 & 表达式 & 适合处理 \\
      \midrule
      代数形式 & \(z=x+\ii y\) & 加法 \\
      三角形式 & \(z=|z|(\cos\theta+\ii\sin\theta)\) & 辐角 \\
      指数形式 & \(z=|z|\ee^{\ii\theta}\) & 乘法、指对幂 \\
      \bottomrule
    \end{tabularx}
  \end{center}
  \item 实部：\(\operatorname{Re}\left ( z \right ) = x\)；虚部：\(\operatorname{Im}\left ( z \right ) = y\)。
  \item 特例：\(y=0\Rightarrow z\) 为实数；$x=0$且$y\neq0\Rightarrow z$ 为纯虚数。
  \item 共轭：\(\overline z=x-y\ii\)。
  \item 运算律：
    \[
      \begin{aligned}
        \overline{z_1\pm z_2}&=\overline{z_1}\pm\overline{z_2},\\
        \overline{z_1z_2}&=\overline{z_1}\,\overline{z_2},\\
        \overline{\left(\frac{z_1}{z_2}\right)}
          &=\frac{\overline{z_1}}{\overline{z_2}}\qquad(z_2\neq0).
      \end{aligned}
    \]
  \item 辐角与其计算

    对 \(z\neq0\)：
    \[
      \arg z\in(-\uppi,\uppi],
      \qquad
      \operatorname{Arg}z=\arg z+2k\uppi,
      \quad k\in\mathbb Z.
    \]

    \[
    \arg z=
    \begin{cases}
      \arctan\dfrac{y}{x}, & x>0,\\[4pt]
      \arctan\dfrac{y}{x}+\uppi, & x<0,\ y>0,\\[4pt]
      \arctan\dfrac{y}{x}-\uppi, & x<0,\ y<0,\\[4pt]
      \dfrac{\uppi}{2}, & x=0,\ y>0,\\[4pt]
      -\dfrac{\uppi}{2}, & x=0,\ y<0,\\[4pt]
      0, & y=0,\ x>0,\\[4pt]
      \uppi, & y=0,\ x<0.
    \end{cases}
    \]
    计算部分其实直接画图算也没什么不好的。
\end{enumerate}

\subsection{复球面}

复分析中称模为无穷大的复数为无穷远点，记为$\infty$。
将$\infty$加入到$\CC$中，便得到了扩充复平面$\CC_\infty$。
表示扩充复平面常用复球面。

\begin{enumerate}
  \item 球极投影

在三维空间中，取单位球面$S^2$，其北极点为$(0,0,1)$，南极点为$(0,0,-1)$。
将复平面$\CC$放在球面$S^2$的赤道平面上，取北极点为投影点，则可将复平面$\CC$上的点
投影到球面$S^2$上。将$\infty$对应到北极点，则得到了复球面$S^2$与扩充复平面$\CC_\infty$的一一对应关系。

  \item 特征：在复球面上，北极点代表$\infty$，南极点为$0$，$1$、$\ii$都在赤道上。
\end{enumerate}

\section{复函数}

\subsection{概念}

设$D\subseteq\CC$，$z\in D$，记$w=f(z)$，则称$f$为定义在$D$上的复函数。
\begin{itemize}
  \item 若每个 \(z\) 对应一个 \(w\)，则为单值函数。
  \item 若每个 \(z\) 对应多个 \(w\)，则为多值函数。
\end{itemize}

若将复函数$f$看作从$D$到$\RR^2$的映射，即令
\[
  z=x+\ii y,
  \qquad
  w=u+\ii v,
\]
则有
\[
  w=f(z)
  \quad\Longleftrightarrow\quad
  (x,y)\longmapsto\bigl(u(x,y),v(x,y)\bigr).
\]

\subsection{初等复函数}
\begin{enumerate}
  \item 常数函数

\[
  w=C\quad(C\in\CC).
\]

  \item 多项式函数

\[
  w=C_0+C_1z+\cdots+C_nz^n\quad(C_i\in\CC)
\]

  \item 指数函数  

\[
  \ee^z=\ee^{x+y\ii}=\ee^x(\cos y+\ii\sin y).
\]

其中$\ee^x$决定模，$y$决定辐角。

  \item 对数函数  

\[
\ln z
=\ln|z|+\ii\operatorname{Arg}z
=\ln|z|+\ii\bigl(\arg z+2k\uppi\bigr),
\qquad k\in\mathbb Z.
\]

$\operatorname{Arg}z$的多值性造就了$\ln z$的多值性。
  
  \item 幂函数  

\[
  w=z^n
\]

在复数域扩展三角函数与双曲函数的定义，便得到了复三角函数与复双曲函数。

  \item 三角函数  

\[
  \sin z=\frac{\ee^{\ii z}-\ee^{-\ii z}}{2\ii},
  \qquad
  \cos z=\frac{\ee^{\ii z}+\ee^{-\ii z}}{2}.
\]

三角函数的许多性质与实数域的三角函数类似，如奇偶性、周期性、零点等，但请注意：
\textbf{复三角函数无界}。

\begin{center}
\begin{tabularx}{0.88\textwidth}{>{\bfseries}p{2.2cm}X X X}
\toprule
函数 & 奇偶性 & 周期性 & 零点 \\
\midrule
\(\sin z\) & 奇 & $2\uppi$ & $n\uppi$ \\
\(\cos z\) & 偶 & $2\uppi$ & $n\uppi + \frac{\uppi}{2}$ \\
\bottomrule
\end{tabularx}
\end{center}

  \item 双曲函数

\[
\sinh z=\frac{\ee^z-\ee^{-z}}{2},
\qquad
\cosh z=\frac{\ee^z+\ee^{-z}}{2}.
\]
\end{enumerate}
\begin{methodbox}[【题型】计算复三角函数]
\[
\sin(a+\ii b)=\sin a\cosh b+\ii\cos a\sinh b.
\]
\[
\cos(a+\ii b)=\cos a\cosh b-\ii\sin a\sinh b.
\]
\end{methodbox}

\section{连续与极限}
复函数的连续与极限的定义与二元实函数类似，在此不多赘述。

\section{微分学}


\subsection{导数}
  
\subsection{微分}

$\dd w = f'(z) \dd z$

\subsection{解析}
  
  \begin{enumerate}
    
    \item 解析：在点$z_0$及其某一邻域均存在导数的复函数称为在点$z_0$解析。
    
    \item 解析函数：在区域$D$内每一点均解析的复函数称为在区域$D$解析。 
  
  \end{enumerate}

  \begin{methodbox}[解析与可导、可微关系]
    
    解析 $\Rightarrow$ 可导 $\Leftrightarrow$ 可微

  \end{methodbox}

\end{document}
